\documentclass[../main.tex]{subfiles}
\begin{document}

\chapter{Variational autoencoders}

\[
\begin{aligned}
\boldsymbol{z} &\sim  p_{\boldsymbol{\theta}}(\boldsymbol{z}) \\
\boldsymbol{x} | \boldsymbol{z} &\sim  \text{Expfam}(\boldsymbol{x} | d_{\boldsymbol{\theta}}(z))
\end{aligned}
\]


where $p(\boldsymbol{z})$ is some kind prior on the latent code $\boldsymbol{z}$, $d_{\boldsymbol{\theta}}(\boldsymbol{z})$ is a deep nn (\emph{decoder}). Posterior inference is computationally intractable, 
\[
p_{\boldsymbol{\theta}}(\boldsymbol{x}) = \int p_{\boldsymbol{\theta}} (\boldsymbol{x} | \boldsymbol{z}) p_{\boldsymbol{\theta}}(\boldsymbol{z}) d \boldsymbol{z}.
\]
Hence we need to approximate inference.
We use \emph{amortized inference}, we train another model $q_{\phi}(\boldsymbol{z} | \boldsymbol{x})$, called the \emph{recognition network}, simultaneously with the generative 
model to approximate posterior inference.

\begin{note}{MNIST Example Note}
   \begin{enumerate}
   \item Prior over latents: we first sample a latent code from a simple prior, typically a standard Gaussian: 
   \[
   \boldsymbol{z} \sim p(\boldsymbol{z}) = \mathcal{N}(0, \boldsymbol{I}).
   \]
   \item Likelihood(decoder): Given $z$, the decoder network $d_{\boldsymbol{\theta}}(z)$ outputs the parameter of a distribution over images. For MNIST, a common choice is a Bernoulli likelihood over pixels:
   \[
   p_{\theta}(\boldsymbol{x} | \boldsymbol{z}) = \text{Bernoulli}(\pi_{\theta}(\boldsymbol{z})), \quad \pi_{\boldsymbol{\theta}}(\boldsymbol{z}) = \sigma (d_{\boldsymbol{\theta}}(z)),
   \]
   where $\pi_{\boldsymbol{\theta}}(\boldsymbol{z}) \in  (0,1) ^{784}$ gives the prob that each pixel is ``on''.
   \end{enumerate} 
   This defines the joint distribution $p_{\theta}(\boldsymbol{x},\boldsymbol{z}) = p(\boldsymbol{z})p_{\boldsymbol{\theta}}(\boldsymbol{x}|\boldsymbol{z})$.

   During training we observe $\boldsymbol{x}$(an image) but we do not observe $z$. To learn good parameters $\theta$, we need to perform inference over the latens $\boldsymbol{z}$
   \[
   p_{\theta}(\boldsymbol{z}, \boldsymbol{x}) = \frac{p_{\boldsymbol{\theta}}(\boldsymbol{x} | \boldsymbol{z})}{p_{\boldsymbol{\theta}}(\boldsymbol{x})}.
   \]
   The difficulty is the normalizing constant, which is a high-dimensional integral involving a nn. To avoid this, we want to approximate the posterior using a \emph{recognition network}
   \[
   q_{\phi}(\boldsymbol{z} |\boldsymbol{x})  \approx p_{\boldsymbol{\theta}}(\boldsymbol{z} | \boldsymbol{x}).
   \]
   For MINIST we often choose
   \[
   q_{\phi}(\boldsymbol{z} | \boldsymbol{x}) = \mathcal{N}(\mu_{\phi}(x) , \text{diag}\left( \sigma_{\phi}(x)^{2} \right) ).
   \]
\end{note} 


\section{VAE basics}

A VAE defines a generative model of the form
\[
p_{\boldsymbol{\theta}}(\boldsymbol{z},\boldsymbol{x}) = p_{\boldsymbol{\theta}}(\boldsymbol{z}) p_{\boldsymbol{\theta}}(\boldsymbol{x} | \boldsymbol{z}).
\]
where $p_{\boldsymbol{\theta}}(\boldsymbol{z})$ is usually a Gaussian, and $p_{\boldsymbol{\theta}}(\boldsymbol{x} | \boldsymbol{z})$ is usually a product of exponential family distributions.

In addition, VAE fits a recognition model
\[
q_{\phi}(\boldsymbol{z} | \boldsymbol{x}) = q \left( \boldsymbol{z} | e_{\boldsymbol{\phi}}(\boldsymbol{x}) \right)  \approx p_{\boldsymbol{\theta}}(\boldsymbol{z} | \boldsymbol{x})
\]
to perform amortized inference. Here $q_{\phi}(\boldsymbol{z} | \boldsymbol{x})$ is usually a Gaussian, with parameters computed by a nn encoder $e_{\boldsymbol{\phi}}(\boldsymbol{x})$:
\[
\begin{aligned}
q_{\phi}(\boldsymbol{z} | \boldsymbol{x}) &= \mathcal{N}\left( \boldsymbol{z} | \boldsymbol{\mu}, \text{diag} \left( \exp(\boldsymbol{\ell}) \right)  \right)  \\
\left( \boldsymbol{\mu}, \boldsymbol{\ell} \right)  &= e_{\phi}(\boldsymbol{x})
\end{aligned}
\]
where $\boldsymbol{\ell} = \log \boldsymbol{\sigma}$. The amortized inference is introduced in Section~\ref{sec:amortized-vi}.

\paragraph{Model fitting}

WE can fit a VAE using \emph{amortized stochastic varational inference}










\printbibliography
\end{document}
